\documentclass{beamer}

\newcommand{\week}{Financial Markets 6 of 8}

\title{Equities 1: Dividend Discount Model}
\subtitle{Mishkin Chapter 6}
\author{Econ 357}
\date{\week}

% Reference the shared preamble
\input{preamble}

% Set the footer -- change 
\setbeamertemplate{footline}{
  \leavevmode%
  \vspace{2ex}
  \hbox{%
    % Left box: Econ 457
    \begin{beamercolorbox}[wd=.4\paperwidth,ht=2.5ex,dp=1ex,left]{author in head/foot}%
      \hspace{1em}Econ 357
    \end{beamercolorbox}%
    % Middle box: Week
    \begin{beamercolorbox}[wd=.2\paperwidth,ht=2.5ex,dp=1ex,center]{date in head/foot}%
      \centering\week
    \end{beamercolorbox}%
    % Right box: Slide numbers
    \begin{beamercolorbox}[wd=.4\paperwidth,ht=2.5ex,dp=1ex,center]{date in head/foot}%
      \hfill\insertframenumber{} 
    \end{beamercolorbox}%
  }%
  \vskip0pt%
}

\begin{document}

\frame{\titlepage}

\begin{frame}
    \frametitle{Outline for Financial Markets}
        \begin{enumerate}
            \item Financial Market Returns
            \item Bonds 1: Discounting, Prices and Yields
            \item Bonds 2: Nominal v. Real Interest Rates, Supply and Demand
            \item Bonds 3: Market for Money
            \item Bonds 4: Risk and Term Structure
            \item \fbox{Equities 1: Dividend Discount Model}
            \item Equities 2: P/E Ratio
            \item Equities 3: Theories of Stock Pricing
            
        \end{enumerate}
\end{frame}


\begin{frame}
    \frametitle{Outline for Today's Lecture}
    \begin{enumerate}
        \item Generic Asset Pricing Equation
        \item The Price of Equities
        \begin{itemize}
            \item Dividends
            \item One period valuation model
            \item Dividend Discount Model (Gordon Growth Model)
            \item Interpretation
        \end{itemize}
        \item Practice Questions
    \end{enumerate}

\end{frame}

\begin{frame}[t]
    \frametitle{1. Generic Asset Pricing Equation}
    \section{Generic Asset Pricing Equation}

    $$
    \text{Value of an Asset} = \sum_{t=1}^{N} \frac{E(\text{Cash Flow}_t)}{(1+r)^t}
    $$

    For bonds:
    \begin{itemize}
        \item $N$ = maturity
        \item Cash flows = coupons + par value at the end
        \item $r$ = Yield to maturity, observable in the market
    \end{itemize}

    For equities:
    \begin{itemize}
        \item $N$ = $\infty$
        \item Which cash flow to use?
        \item $r$ = discount rate
        \begin{itemize}
            \item Discount rate reflects the riskiness of the investment and investor preferences
            \item Will mostly likely be higher than bond yields
        \end{itemize}
    \end{itemize}
\end{frame}

\begin{frame}[t]
    \frametitle{1. Generic Asset Pricing Equation}

    $$
    \text{Value of an Asset} = \sum_{t=1}^{N} \frac{E(\text{Cash Flow}_t)}{(1+r)^t}
    $$

    Notes:
    \begin{itemize}
        \item Note that the left hand side says “value.”   This may be different for different people, depending on their discount rate and their expectations (more on this later). 
        \item More importantly, this is not necessarily equal to the price.   In fact, your value of an asset may differ from the price for a number of reasons.
    \end{itemize}


\end{frame}

\begin{frame}[t]
    \frametitle{1. Generic Asset Pricing Equation}

    Notes on notation:
    \begin{itemize}
        \item Often use $k$ as the discount rate for equities.
        \item Income for equities will be the dividends, indicated with a $D$
    \end{itemize}

\end{frame}

\begin{frame}[t]
    \frametitle{2. The Price of Equities}
    \framesubtitle{Dividends}
    \section{The Price of Equities}
    \subsection{Dividends}

    \begin{itemize}
        \item Discretionary distribution of earnings to shareholders.   Management could alternatively decide to reinvest in the company.
        \item Usually paid quarterly
        \item NOT guaranteed.   For example:
        \begin{itemize}
            \item Apple stopped paying a dividend after the 1990s tech bubble, resumed in 2012
            \item Many banks cut or suspended dividends in 2008; regulators prevents banks from increasing dividends in 2020
        \end{itemize}
        \item Recently companies have preferred buying back stock as a way to distribute cash to shareholders.
    \end{itemize}

\end{frame}

\begin{frame}[t]
    \frametitle{2. The Price of Equities}
    \framesubtitle{Dividends}

    \vspace{-1em}
    $$
    \text{Annualized Dividend Yield} = \frac{(4 \times \text{Quarterly Dividend} (\$))}{(\text{Price of Stock} (\$))}
    $$
    \vspace{1em}

    % Requires: \usepackage{booktabs}
    \begin{table}
    \centering
    \caption{\textbf{Dividend Data for Selected Companies}}
    \footnotesize
        \begin{tblr}{
      colspec = {Q[l,wd=4cm] Q[c,wd=2cm] Q[c,wd=2cm]}
      }
            \toprule
            Company & Dividend & Dividend Yield \\
            \midrule
            Amazon & \$0 & \\
            Microsoft & \$0.91 & 0.9\% \\
            Wells Fargo & \$0.45 & 1.92\% \\
            Coca-Cola & \$0.51 & 2.60\% \\
            Ford Motor & \$0.15 & 4.35\% \\
            Edison International & \$0.875 & 5.48\% \\
            Verizon & \$0.71 & 6.10\% \\
            \bottomrule
    \end{tblr}
    \end{table}

    \blfootnote{Source: Google Finance, as of Feb 2026}

\end{frame}

\begin{frame}[t]
    \frametitle{2. The Price of Equities}
    \framesubtitle{Dividends}

    Why use dividends?
    \begin{itemize}
        \item Starting point
        \item Can augment with an expected price at end of holding period
    \end{itemize}
    \vspace{2em}

    Why \textbf{not} use dividends?
    \begin{itemize}
        \item Out-dated.   Companies purchase stock with cash.
        \item What about high growth companies?
    \end{itemize}

\end{frame}

\begin{frame}[t]
    \frametitle{2. The Price of Equities}
    \framesubtitle{One period valuation model}
    \subsection{One period valuation model}

    $$
    V_0 = \frac{D}{(1+k)} + \frac{E(P_1)}{(1+k)}
    $$

    Where:
    \begin{itemize}
        \item $V_0$ = current value ($t=0$)
        \item $E(P_1)$ = expected price at the end of the period
        \item $D$ = dividend (income you expect to receive during the period)
        \item $k$ = required rate of return
    \end{itemize}

\end{frame}

\begin{frame}[t]
    \frametitle{2. The Price of Equities}
    \framesubtitle{One period valuation model}

    $$
    V_0 = \frac{D}{(1+k)} + \frac{E(P_1)}{(1+k)}
    $$
    \vspace{1em}

    \textit{Example Problem:}\\
    
    Your discount rate for stock ABC is 12\%.   You expect stock ABC to pay a dividend of \$0.16 and you forecast the stock price will be \$60 a year from now.   What do you think the price should be today?   You notice that the stock is priced at \$50 today.   What should you do?

    \blfootnote{Mishkin refers to the discount rate in the following way: “you would be satisfied with earning a 12\% return”}

\end{frame}

\begin{frame}[t]
    \frametitle{2. The Price of Equities}
    \framesubtitle{Dividend Discount Model}
    \subsection{Dividend Discount Model}

    $$
    V_0 = \frac{D_1}{(1+k)} + \frac{D_2}{(1+k)^2} + \frac{D_3}{(1+k)^3} + \dots$$
    \vspace{1em}

    Where:
    \begin{itemize}
        \item $V_0$ = current value
        \item $D_t$ = dividend expected at time $t$
        \item $k$ = required rate of return (discount rate)
    \end{itemize}
    \vspace{1em}
    
    In words: The value of an equity equals the present discounted value of all expected future dividends

\end{frame}

\begin{frame}[t]
    \frametitle{2. The Price of Equities}
    \framesubtitle{Gordon Growth Model}

    $$
    V_0 = \frac{D_1}{(1+k)} + \frac{D_2}{(1+k)^2} + \frac{D_3}{(1+k)^3} + \dots$$
    \vspace{1em}
    
    Assume that the growth rate of equities is constant at rate $g$, then this equation can be written:

        $$
    V_0 = \frac{D_0(1+g)}{(1+k)} + \frac{D_0(1+g)^2}{(1+k)^2} + \frac{D_0(1+g)^3}{(1+k)^3} + \dots$$

\end{frame}

\begin{frame}[t]
    \frametitle{2. The Price of Equities}
    \framesubtitle{Gordon Growth Model}

    Further assume that the growth rate is lower than the discount rate ($g<k$), then this equation can be simplified to:
    \vspace{1em}

    $$
    V_0 = \frac{D_0(1+g)}{(k-g)}
    $$


    This equation is known as the \textbf{Gordon Growth Model.}   Note that it requires two assumptions: constant growth rate ($g$) and $g<k$.

    \blfootnote{The derivation of this is not hard, will be presented in class}
    
\end{frame}

\begin{frame}[t]
    \frametitle{2. The Price of Equities}
    \framesubtitle{Gordon Growth Model: Interpretation}
    \subsection{Interpretation}

    The Gordon Growth Model implies that a stock’s value will be greater:
    \begin{enumerate}
        \item The larger its expected dividend per share.
        \item The lower the market capitalization rate, k.
        \item The higher the expected growth rate of dividends.
    \end{enumerate}

\end{frame}

\begin{frame}[t]
    \frametitle{2. The Price of Equities}
    \framesubtitle{Gordon Growth Model: Interpretation}

    The stock price is expected to grow at the same rate as dividends ($g$).   (it is easy to show this, the growth rate is $\frac{V_1}{V_0}$, write this out and see everything that cancels)
    \vspace{1em}
    
    Another implication is that we can rewrite the expected return of a stock as dividend yield + growth rate of dividends:
    
    \begin{align*}
    E(r) &= \text{Dividend Yield} + \text{Capital gains yield} \\
    &= \frac{D_1}{P_0} + \frac{P_1 - P_0}{P_0} \\
    &= \frac{D_1}{P_0} + g
    \end{align*}
    
\end{frame}

\begin{frame}[t]
    \frametitle{2. The Price of Equities}
    \framesubtitle{Gordon Growth Model: Extensions}
    \subsection{Extensions}

    Possible Extensions:
    \begin{enumerate}
        \item High-growth companies could be modeled in two stages: fast growth in the first stage, lower and stable growth in the second stage.
        \item Could use cash flows instead of dividends for companies that don’t pay dividends
    \end{enumerate}
    
\end{frame}

\begin{practiceframe}[t]
    \frametitle{3. Practice}

\begin{enumerate}
    \item Compute the price of a share of stock that pays \$1 per year dividend and that you expect to be able to sell in one for \$20, assuming you require a 15\% return.
    \item After careful analysis, you have determined that a firm’s dividends should grow at 7\%, on average, in the foreseeable future.   The firm’s last dividend was \$3.  Compute the current price of the stock, assuming the required return is 18\%.
    \item The current price of a stock is \$65.88.   If dividends are expected to be \$1 per share for the next five years, and the required return is 10\%, what should be the price of the stock be in 5 years when you plan to sell it?   If the dividend and the required return are the same, and the stock price is expected to increase by \$1 five year from now, does the current stock price also increase by \$1?  Why or why not?
\end{enumerate}




\end{practiceframe}

\end{document}